\section{Interpretation}

\subsection{Review is not equivalent to validation}

The model distinguishes manuscript scrutiny from scientific validation. Reviewers can improve clarity, identify methodological weaknesses, or prevent some false claims from receiving early legitimacy. They generally do not independently reproduce experiments, regenerate data, verify every derivation, or observe future predictive success. Those activities occur downstream through continuing evidence.

A process can therefore be useful without being constitutive of the scientific method. The relevant institutional question is not whether reviewers sometimes help. It is whether granting a small preliminary committee control over publication attention produces a positive net effect after false positives, false negatives, delay, labor, bias, and lost exploration are all counted.

\subsection{Meaning of the game-theoretic harm claim}

The statement that peer review always harms science is not a claim that every review report causes damage. It is a comparative institutional claim. Holding expert evaluation, criticism, repair, replication, and continuing scrutiny fixed, adding a mandatory prepublication veto introduces exclusion, delay, prestige amplification, strategic conformity, and false-negative lock-in. Those mechanisms can move the scientific system toward a worse equilibrium than the same evaluative process operating without publication control.

The simulations test this system-level claim across broad realistic, symmetric, high-consequence, and equal-labor parameter families. They strongly support the proposed mechanism in those tested regimes. They do not yet constitute a mathematical proof over every possible parameterization. A literal universal theorem would require explicit assumptions under which the no-veto institution weakly dominates the otherwise equivalent gatekeeping institution.

\subsection{The asymmetry of false negatives}

False positives and false negatives are not necessarily equal. The harm associated with one false paper is usually bounded by later correction, although high-stakes misinformation can be severe. A false negative can be irreversible if the work is abandoned, deprived of funding, not replicated, or later rediscovered without attribution. When scientific value is heavy-tailed, one suppressed high-value discovery can outweigh many successful routine filters.

\subsection{Attention, not storage, is scarce}

Digital repositories solve physical distribution, but they do not solve scientific attention. In a literature containing hundreds of millions of records, no researcher can examine every relevant claim. Citation counts, journal reputation, institutional affiliation, recommendation systems, and existing prestige become attention-allocation mechanisms. These signals can create cumulative advantage:

\begin{equation}
\text{attention}\rightarrow\text{citations}\rightarrow\text{credibility}
\rightarrow\text{more attention}.
\end{equation}

An open system without triage can recreate the same concentration problem. This is why unfiltered open publication did not consistently win. The best-performing architecture was generally immediate publication combined with transparent triage and deliberate exploration of neglected work.

\section{Limitations}

The simulations are structural counterfactuals, not empirical estimates of historical peer review. Several limitations are fundamental.

\begin{enumerate}
    \item \textbf{Parameter calibration.} The reviewer-behavior submodel is now constrained to selected controlled-study targets for error detection, disagreement, and outcome bias. The complete ecosystem is not empirically identified: attention loss after rejection, abandonment, resubmission, discoverability, social-value distributions, and comparative open-triage performance still require direct measurement.
    \item \textbf{Truth representation.} Truth is represented by three latent states. Real claims can be valid only within narrow domains, depend on auxiliary assumptions, or change meaning as methods improve.
    \item \textbf{Evidence model.} Evidence is summarized by noisy scalar signals. Real evidence varies in design quality, dependence, incentives, replication fidelity, and susceptibility to coordinated error.
    \item \textbf{Institutional adaptation.} Authors, reviewers, journals, and platforms would adapt strategically to any system. The present model includes some incentives but not full evolutionary adaptation.
    \item \textbf{Social cost weights.} High-consequence weights are illustrative. Different ethical frameworks will value false positives, false negatives, delay, and equity differently.
    \item \textbf{Open-system governance.} Transparent triage is represented abstractly. A real platform would require identity safeguards, conflict-of-interest rules, data standards, fraud controls, auditability, and resistance to manipulation.
\end{enumerate}

The simulations do not test the removal of expert evaluation, and nothing in the model supports that interpretation. They test whether adding a publication veto improves outcomes when expert criticism and continuing evaluation are otherwise retained. Across the tested regimes, the veto produces the game-theoretic harm mechanism described above. Because some deliberately idealized sampled worlds allow peer review to win, the current simulation evidence is not yet a universal dominance proof. The empirical calibration layer strengthens the reviewer-behavior component, but it does not convert unmeasured ecosystem parameters into observations or make the institutional win shares empirical probabilities.

\section{Empirical Tests Suggested by the Model}

The model identifies quantities that can be measured directly:

\begin{enumerate}
    \item the probability that a rejected paper receives equivalent later attention after controlling for intrinsic quality;
    \item reviewer error-detection rates under controlled planted-error designs;
    \item inter-reviewer agreement and decision reversal under duplicate committees;
    \item prestige and institutional-affiliation effects under randomized identity masking;
    \item the proportion of rejected papers later accepted elsewhere and their long-run replication, correction, and predictive outcomes;
    \item changes between preprints and final journal versions, classified as presentational, methodological, statistical, or truth-relevant;
    \item the downstream social cost of delayed high-consequence findings;
    \item recognition gaps for equivalent work from differently resourced authors;
    \item performance of immediate-publication platforms using randomized audit and replication allocation.
\end{enumerate}

A decisive real-world test would randomize comparable manuscripts among traditional review, immediate open publication, and hybrid publication, allocate equal evaluation resources, and track calibration, replication, correction, and value recovery over a long horizon.

\section{Conclusion}

Formal peer review can improve manuscripts and prevent some defective claims from receiving early legitimacy. Those local benefits do not establish that mandatory prepublication gatekeeping maximizes scientific truth recovery.

In the canonical simulation, peer review became competitive only when reviewers were highly accurate, reviewer noise was low, prestige and novelty biases were absent, rejected work retained nearly equal attention, resubmission was easy, and delay and labor were treated as costless. The robustness extensions strengthened the comparison. When every system received the same defect-detection and manuscript-repair mechanisms, peer review won 0.4\% of worlds. When total lifetime expert labor was made exactly equal and the remaining assumptions deliberately favored peer review, it won 14.0\% of worlds, compared with 60.1\% for open triage. Under more realistic structural assumptions, immediate publication with continuing transparent triage recovered more true scientific value and produced lower calibration error. When the asymmetric social cost of missed cancer, public-health, safety, climate, and infrastructure research was included, the penalty for false negatives became substantial.

The new reviewer calibration reproduces controlled findings on planted-error detection, reviewer disagreement, and positive-outcome bias without equating those observed percentages with latent score coefficients. This narrows one major source of arbitrariness. The institutional conclusion nevertheless remains conditional because the downstream consequences of rejection are not yet identified.

The central conclusion is conditional but important:

\begin{quote}
If uncertain prepublication decisions materially reduce later attention, mandatory gatekeeping can impede long-run scientific truth recovery even while successfully filtering some false claims.
\end{quote}

Peer review should therefore be evaluated as an institutional intervention, not assumed to be synonymous with scientific scrutiny or scientific validity. The burden is empirical: compare it with alternatives under equal resources and measure both what it admits and what it prevents from being tested.

\section*{Code and Reproducibility}

The complete Python implementation, legacy baseline, robustness runner, run instructions, and compact result summaries are distributed with this manuscript. The canonical entry point is:

\begin{verbatim}
python src/scientific_publication_simulator.py \
  --worlds 5000 --papers 300 --periods 30 \
  --output-dir simulation_output
\end{verbatim}

The two final robustness experiments can be reproduced with:

\begin{verbatim}
python src/reproduce_robustness_experiments.py \
  --output-root reproduced_results
\end{verbatim}

The empirical reviewer calibration can be regenerated with:

\begin{verbatim}
python src/calibrate_empirical_model.py
\end{verbatim}

The calibration target metadata are stored in \texttt{data/empirical\_targets.json}, and parameter classifications are documented in \texttt{ASSUMPTIONS.md}. The robustness command regenerates the 125,000-paper symmetric correction tables and the 176,000-paper equal-total-labor tables. The simulator requires Python 3, NumPy, and pandas. Fixed random seeds are used for exact reproduction of the committed summaries. Large-scale results should also be rerun across multiple seeds and independently audited before inferential claims are made.
